\chapter{Algorithm Design and Data Structures}\label{ch:algo}

\section{Identification of Critical Operations}
Most of the system is straightforward reading and writing of rows. A few operations,
however, decide whether the system is \emph{correct}, and they are where algorithm design
matters. Table~\ref{tab:critical} lists them.

\begin{table}[H]
\centering\small
\caption{Critical operations and what makes them critical}\label{tab:critical}
\begin{tabularx}{\linewidth}{@{}p{0.9cm}p{4.2cm}X@{}}
\toprule \textbf{ID} & \textbf{Operation} & \textbf{Why it is critical}\\ \midrule
A1 & Resolve and check permissions on every request & Security decision made hundreds of times per minute; must be cheap and always current\\
A2 & Book an appointment slot & Two users may compete for the last seat; overbooking is a visible failure\\
A3 & Approve an ownership transfer & A vehicle must have exactly one current owner before and after\\
A4 & Record a payment attempt & A retried request must not charge twice or lose a failed attempt\\
A5 & Validate a status transition & Wrong moves (for example reopening a completed application) corrupt history\\
A6 & Build a filtered, sorted, paged list & Used by every list screen; must be safe against injection and bounded in size\\
A7 & Prevent a cycle in the region hierarchy & A region that is its own ancestor would break tree traversal\\
A8 & Order the creation of 63 tables & A table cannot be created before the tables it references\\
\bottomrule
\end{tabularx}
\end{table}

\section{Algorithm Design}
The algorithms are given in pseudocode that follows the structure of the implementation.
Lines that start with \textsf{LOCK} correspond to \cmd{SELECT ... FOR UPDATE}.

\subsection{A1: permission resolution and check}
\begin{lstlisting}[language={},mathescape=false]
resolveUser(userId):
    user  <- select user by id; if none or inactive: 401
    roles <- select role names via user_roles
    perms <- SET of distinct permission keys via
             user_roles JOIN role_permissions JOIN permissions
    emp   <- active employee of user's person (if any)
    return CurrentUser(user, roles, perms, emp, ...)

check(handler, user):
    if handler is @Public: allow
    if user is none: 401
    for key in @Requires(handler): if key not in user.perms: 403
    if @RequiresAny(handler) and none of its keys in user.perms: 403
    allow
\end{lstlisting}
The permission set is a hash set, so each membership test is $O(1)$. Loading it costs a
few indexed joins per request; because it is read each time, a change to a role takes
effect immediately.

\subsection{A2: booking a slot}
\begin{lstlisting}[language={},mathescape=false]
book(application, slot):
    LOCK application
    LOCK slots in ascending id order        // TreeSet of ids
    require slot.office == application.office
    require slot start is in the future
    n <- UPDATE slot SET booked = booked + 1
         WHERE id = slot AND booked < capacity   // guarded update
    if n == 0: 409 SLOT_FULL
    token <- "S{slot}-{ (highest suffix among slot tokens) + 1 }"
    insert appointment(BOOKED, token); insert audit row
    queue notification (delivered after commit)
\end{lstlisting}
Three layers prevent overbooking. The lock serialises competing bookers; the guarded
\cmd{UPDATE ... WHERE booked_count < capacity} is an atomic compare-and-set that refuses
to increment past the capacity even if the lock were somehow missed; and the table's
CHECK constraint (\cd{booked_count <= capacity}) is the last line of defence. The token
is computed while the slot is locked, so two callers can never receive the same number.
Figure~\ref{fig:flowbook} shows the flow.

\fig{flow_booking.png}{Flowchart of slot booking}{fig:flowbook}{0.72}

\subsection{A3: approving an ownership transfer}
\begin{lstlisting}[language={},mathescape=false]
approve(transferId, date):
    LOCK transfer;  require status == PENDING
    LOCK vehicle;   require status == ACTIVE
    LOCK current ownership row (effective_to IS NULL)
    require ownership.citizen == transfer.from_citizen
    require transferee not blacklisted
    require ownership.effective_from <= date <= today
    ownership.effective_to <- date       // close the old row
    FLUSH                                 // old row must reach the DB first
    insert ownership(vehicle, to_citizen, effective_from = date)
    transfer.status <- APPROVED; transfer.approved_at <- now
    insert audit rows; queue notifications
    COMMIT   (any failure above: ROLLBACK everything)
\end{lstlisting}
The explicit flush matters. The unique key on the generated column
\cd{vehicle_ownerships.current_flag} allows one open row per vehicle, so the old row must be
closed in the database \emph{before} the new open row is inserted; otherwise the insert would
violate the key. A dedicated test (\texttt{aFailureMidApprovalRollsEverythingBack}) forces
a failure after the first changes and checks that nothing remains. Figure~\ref{fig:flowtransfer}
shows the flow.

\fig{flow_transfer.png}{Flowchart of ownership-transfer approval}{fig:flowtransfer}{0.75}

\subsection{A4: idempotent payment}
\begin{lstlisting}[language={},mathescape=false]
payTarget(type, id, reference, outcome):
    prior <- attempt with this gateway reference          // indexed lookup
    if prior exists:
        if prior.payment is for the same (type, id): return it   // replay
        else: 409 GATEWAY_REFERENCE_REUSED
    LOCK target;  if a SUCCESS payment exists: 409 ALREADY_PAID
    payment <- reuse open PENDING/FAILED payment or create one
               (amount = amount due, never taken on trust)
    LOCK payment
    insert attempt(number = max + 1, reference, outcome)
    SUCCESS -> payment SUCCESS, settle target (challan PAID ...)
    FAILED  -> payment FAILED (retry allowed)
    TIMEOUT -> payment stays PENDING (result unknown)
\end{lstlisting}
The gateway reference acts as an \emph{idempotency key}: the same reference always yields
the same result, and the replay is answered \emph{before} the target is validated, because
the item may legitimately be paid by now. The \cd{idx_attempt_gateway_ref} index keeps this
lookup fast (Chapter~\ref{ch:perf}). A timeout deliberately does not mark the payment failed:
the gateway outcome is unknown, so the payment stays pending and can be reconciled.
Figure~\ref{fig:flowpay} shows the flow.

\fig{flow_payment.png}{Flowchart of an idempotent payment}{fig:flowpay}{0.75}

\subsection{A5: status transition validation}
Each workflow has a table of legal moves, a map from a status to the set of statuses
that may follow it. A move is accepted only if the target is in the set of the current
status, and some targets additionally require a permission and a reason.
\begin{lstlisting}[language={},mathescape=false]
TRANSITIONS = { SUBMITTED: {DOCS_PENDING, UNDER_VERIFICATION, REJECTED, CANCELLED},
                ... , APPROVED: {COMPLETED}, REJECTED: {}, COMPLETED: {}, CANCELLED: {} }

applyStatus(app, next, reason):
    if next not in TRANSITIONS[app.status]: 409 INVALID_TRANSITION
    if next in {REJECTED, CANCELLED, DOCS_PENDING} and reason is empty: reject (validation error)
    require permission PERMISSION_FOR_TARGET[next]
    append history row(prev, next, user, reason); app.status <- next
\end{lstlisting}

\subsection{A6: filtered, sorted, paged lists}
A small query builder appends a condition only when its value is present, binds every value
as a parameter, escapes the characters \texttt{\%} and \texttt{\_} in search text, and
accepts a sort key only if it is in a whitelist. A list request runs a \texttt{COUNT} over
the filtered query and then one bounded page (at most 100 rows), and returns the items with
the page number, page size, total and page count.

\subsection{A7: region hierarchy}
Regions form a tree through \cd{parent_region_id}. Before a parent is assigned, the
service walks up the parent chain from the proposed parent; if it meets the region
itself, the change is refused with \cd{REGION_CYCLE}. The region tree for display is built in
one pass: every region becomes a node in a hash map by id, and each node with a parent is
appended to its parent's list of children.

\subsection{A8: dependency-ordered schema creation}
Treat the 63 tables as vertices and each foreign key as a directed edge from the
referencing table to the referenced table. A table can be created only after all the tables
it points to, so a valid creation order is a topological order of this directed acyclic
graph. We computed the longest-path level of every table and checked that the order used
by the scripts respects every edge. There are $V=63$ vertices and $E=90$ edges once the
single self-reference (\cd{regions}) is set aside, and the levels run from 0 to 6
(Table~\ref{tab:levels}). Topological ordering and level computation take $O(V+E)$ time,
which here is 153 elementary steps; the check found no edge pointing backwards.

\begin{table}[H]
\centering\small
\caption{Dependency depth of the 63 tables}\label{tab:levels}
\begin{tabularx}{\linewidth}{@{}cp{1.6cm}X@{}}
\toprule \textbf{Level} & \textbf{Tables} & \textbf{Examples}\\ \midrule
0 & 18 & document\_types, fuel\_types, departments, designations, challans, audit\_logs\\
1 & 12 & citizens, employees, addresses, fee\_structures, notifications\\
2 & 11 & applicants, appeals, complaints, counters, driving\_schools, employee\_postings\\
3 & 9 & applications, appointment\_slots, driving\_instructors, insurance\_policies\\
4 & 9 & documents, appointments, driving\_tests, fitness\_certificates, challan\_violations\\
5 & 2 & driving\_licences, permit\_status\_history\\
6 & 2 & licence\_class\_assignments, licence\_status\_history\\
\bottomrule
\end{tabularx}
\end{table}

\section{Flowcharts}
The flowcharts of the three transactional algorithms are Figures~\ref{fig:flowbook},
\ref{fig:flowtransfer} and~\ref{fig:flowpay}; the permission decision is
Figure~\ref{fig:perm}; the request life cycle is Figure~\ref{fig:request}.

\section{Data Structure Selection}
Table~\ref{tab:ds} lists the data structures that matter and why each was chosen.

\begin{table}[H]
\centering\small
\caption{Data structures and where they are used}\label{tab:ds}
\begin{tabularx}{\linewidth}{@{}p{3.1cm}p{4.3cm}X@{}}
\toprule \textbf{Structure} & \textbf{Used for} & \textbf{Reason}\\ \midrule
Hash map of sets & Status transition tables for applications, licences, permits, vehicles, challans, complaints and appeals & Constant-time lookup of the allowed next states\\
Hash set & A user's permission keys; redacted audit keys; allowed file extensions & Constant-time membership tests\\
Sorted set (tree set) & Slot ids to lock, so that locks are always taken in ascending order & Gives a canonical order and removes duplicates in $O(s\log s)$\\
Concurrent hash map & Revoked token ids with their expiry times & Safe for many threads; expired entries are purged when a new one is added\\
Insertion-ordered map & Audit values and JSON bodies & Keeps field order stable in the log\\
Tree (adjacency by parent id) & Region hierarchy & Natural for ``belongs to'' with unlimited depth\\
Directed acyclic graph & Foreign-key dependencies & Defines a valid table creation order\\
B$^+$-tree index (InnoDB) & Every primary, unique and secondary key & $O(\log n)$ point and range lookups\\
Table used as a queue & Notifications: QUEUED then SENT or FAILED & Durable before delivery is attempted; survives a restart\\
List & Result pages, sorted permission lists, collected rows & Simple ordered collections\\
\bottomrule
\end{tabularx}
\end{table}

\section{Algorithm Selection and Justification}
\paragraph{Pessimistic locking over optimistic versioning.} The contested rows (the last
seat of a slot, the single current owner, the payment of one challan) are exactly the
rows where conflicts are likely. Optimistic locking would let both writers do the work
and make one of them retry, which is wasteful for hot rows. A short row lock is cheaper
and keeps the code simple. The cost is that competing writers wait, which is acceptable
for office-scale traffic.

\paragraph{Ordered locking against deadlock.} A deadlock needs a cycle in the
``waits-for'' relation between transactions. If every transaction acquires locks of the
same kind in a single global order, such a cycle is impossible, because a transaction
waiting for a lock of higher rank can never be waited for by the holder of that lock
unless it also held a lock of higher rank. The implementation fixes the order inside each
operation family: application, then appointment, then slots in ascending id; transfer,
then vehicle, then ownership row; target, then payment. When two slots are involved, as in
a reschedule, the ids are sorted first. The database can still detect a rare deadlock and
abort one victim; the error mapper turns that into a ``please retry'' response.

\paragraph{Guarded update over read-modify-write.} Reading \texttt{booked\_count},
adding one in Java and writing it back is a race. Putting the condition in the
\texttt{WHERE} clause makes the check and the change a single atomic statement.

\paragraph{Idempotency key over de-duplication by time.} Rejecting a ``similar payment a few
seconds ago'' would be a guess. A unique reference supplied by the caller gives an exact
rule.

\paragraph{Offset paging.} Lists use \texttt{LIMIT}/\texttt{OFFSET} because office lists are
short and users jump to arbitrary pages. For very large tables, key-set paging would be
faster; this is discussed in Section~\ref{sec:opt}.

\section{Time Complexity Analysis}
Let $n$ be the number of rows in a table, $k$ the number of rows matched by a query, $p$
the page size, $s$ the number of slots locked together, $b$ the number of bookings in one
slot and $d$ the depth of the region tree.

\begin{table}[H]
\centering\small
\caption{Time complexity of the critical operations}\label{tab:time}
\begin{tabularx}{\linewidth}{@{}p{4.3cm}p{2.9cm}X@{}}
\toprule \textbf{Operation} & \textbf{Cost} & \textbf{Explanation}\\ \midrule
Permission membership test & $O(1)$ expected & Hash set lookup\\
Load a user's permissions & $O(\log n + k)$ per query & Indexed joins over the user's role and permission rows\\
Status transition check & $O(1)$ expected & Map lookup, then set membership\\
Lock one row by key & $O(\log n)$ & Primary-key index descent\\
Slot booking & $O(s\log s + s\log n + b)$ & Sort ids, lock $s$ rows, guarded update, scan the slot's $b$ tokens for the next number\\
Ownership approval & $O(\log n)$ & A constant number of indexed lookups and writes\\
Payment replay lookup & $O(\log n)$ with the index; $O(n)$ without & Lookup by gateway reference\\
Paged list & $O(\log n + p)$ to return a page; offset adds $O(\text{offset})$ & Count plus bounded page\\
Region cycle check & $O(d)$ queries & Walk up the chain\\
Region tree build & $O(n)$ & One pass with a hash map\\
Schema ordering and check & $O(V+E)$ & Topological order of the foreign-key graph\\
\bottomrule
\end{tabularx}
\end{table}
The difference between $O(n)$ and $O(\log n)$ for the payment lookup is not only
theoretical: on 500\,000 rows, the lookup took 405.9~ms as a full scan and 1.4~ms through
the index (Section~\ref{sec:indexres}).

\section{Space Complexity Analysis}
The algorithms need little extra memory. A permission set holds at most 62 strings per
request. Locking $s$ slots keeps $s$ entities in memory, $O(s)$. A page holds at most $p\le
100$ rows. The region tree holds $O(n)$ nodes for $n$ regions. The revoked-token map holds one
entry per revoked token until its expiry, so its size is proportional to the number of
sign-outs and refreshes within one token lifetime. The topological check needs $O(V+E)$
space for the adjacency data. Persistent storage grows with the number of records and is
governed by the database.

\section{Optimization Opportunities}\label{sec:opt}
\begin{itemize}
\item \textbf{Key-set paging.} For tables with millions of rows, replace the offset by a ``rows after this key'' condition so that deep pages cost $O(\log n + p)$.
\item \textbf{Cache permissions briefly.} A per-user cache with a short lifetime would remove the two permission queries per request, at the price of a small delay before a role change takes effect; we chose freshness over speed.
\item \textbf{Token sequence per slot.} The next token number is found by scanning the slot's tokens; a counter column on the slot would make it $O(1)$.
\item \textbf{Shared revocation store.} Replace the in-memory revoked-token map by a shared store before running several backend instances.
\item \textbf{Batch notification delivery.} Deliver queued messages in batches by a scheduled job instead of one transaction per message when volume grows.
\item \textbf{Measure first.} Each of these should be justified by a measurement; the performance chapter shows where the time actually goes.
\end{itemize}
